\section{The tree of global cut points}\label{sec:pretree}

We now isolate the part of the boundary argument that uses its global cut
points.  Throughout this section, $B$ is a nondegenerate compact metric
continuum on which a finitely presented, torsion-free, one-ended group
$\Gamma$ acts minimally as a convergence group.  We also assume that
$\Gamma$ has no nontrivial splitting over a finite or two-ended subgroup.
These properties have been established for the boundary constructed above.
Our aim is to place the global cut points in a simplicial tree with finite
quotient, retaining which cut points separate which. In the next section,
we will project onto the cut points adjacent to one of the tree's other
vertices. The closure of that neighboring set will be the continuum
without cut points needed for the rest of the proof.

\subsection{The cut-point pretree}

A point $p\in B$ is a \emph{global cut point} if
$B\setminus\{p\}$ is disconnected.  Write $P$ for their set.  We shall only
need the present construction when $P\ne\varnothing$.  For distinct
$a,p,b\in P$, write $apb$ when an open separation of
$B\setminus\{p\}$ places $a$ and $b$ on opposite sides.

This relation is a \emph{pretree}: it is a strict betweenness relation with
the following properties:
\[
 \neg(aba),\qquad
 abc\Longleftrightarrow cba,\qquad
 \neg(abc\wedge acb),\qquad
 abc, d\ne b\Longrightarrow abd\ \text{or}\ cbd.
\]
Write $(a,b)=\{p:apb\}$ and $[a,b]=(a,b)\cup\{a,b\}$, with
$[a,a]=\{a\}$.  A subset is \emph{full} if it contains $[a,b]$ whenever
it contains $a,b$.  A pretree is \emph{median} if every triple has a point
in the three pairwise closed intervals; this point is unique and is
denoted by $\operatorname{med}(a,b,c)$.  It is \emph{complete} if every
nonempty full linearly ordered subset is an interval, allowing open and
half-open intervals.  These are order-theoretic notions; we do not assert
that intervals in the cut-point pretree are arcs in $B$.

We use the following precise part of Bowditch's construction.  It supplies
the established pretree machinery; the argument after it proves the
additional assertion needed when the cut points are parabolic.

\begin{proposition}[Bowditch's cut-point construction]\label{pret:package}
Let a group act minimally by convergence on a nondegenerate compact metric
continuum $B$, and let $P\ne\varnothing$ be its set of global cut points.
There is a complete median pretree $\Phi$ containing $P$ equivariantly,
with the same betweenness on $P$, with these properties.
\begin{enumerate}[label=\textup{(\roman*)}]
\item\label{pret:item-adjacency}
Two distinct points are called adjacent if their open interval is empty.
Every adjacent pair in $\Phi$ has one point in $P$ and one outside $P$.
A full subset whose intervals are finite is the vertex set of a simplicial
tree, with these adjacent pairs as edges.
\item\label{pret:item-support}
There is an equivariant map $\phi:B\to\Phi$, equal to the identity on
$P$, and, for every $q\in\Phi$, an equivariantly assigned nonempty compact
subset $R_B(q)\subset B$.  For $q\in P$ this set is $\{q\}$; for
$q\notin P$ it satisfies
\begin{equation}\label{pret:support}
 R_B(q)=\phi^{-1}(q)\cup\Lambda(q),
 \qquad \Lambda(q)=\{p\in P:p\text{ is adjacent to }q\}.
\end{equation}
No point of $\Phi$ is fixed by the whole group.
\item\label{pret:item-quotient}
Starting with equality on $\Phi$, successively collapse finite-interval
classes (two points belong to one such class when their closed interval
in the current quotient is finite), and at limit ordinals take the union
of the preceding equivalence
relations.  This process stabilizes.  Its final quotient is the canonical
dendritic quotient $D(B)$ of $B$.  In particular, $D(B)$ is a dendrite and
there is a continuous equivariant surjection $B\to D(B)$.  The induced
action on $D(B)$ is minimal and convergence.
\item\label{pret:item-fixed}
If a loxodromic element fixes $q,r\in\Phi$ with $[q,r]$ infinite, then
$q,r$ are terminal in $\Phi$: neither lies strictly between two points.
If a parabolic element has its unique fixed point in $B\setminus P$,
then it has exactly one fixed point in $\Phi$.
\end{enumerate}
\end{proposition}

\begin{proof}[References for the construction]
The complete median pretree $\Phi$ is the flow completion of $P$ in
\cite[Theorem~3.19]{Bowditch1999Treelike}.  The embedding preserves the
separation betweenness of the cut set.  Adjacency and the simplicial-tree
description are Lemmas~3.28 and~3.34.  The compact supports are constructed
in Section~5; their formula is Lemma~6.8, and the absence of a fixed point
is Lemma~6.11.  Full equivalence relations and their successive
finite-interval collapses are treated in Lemmas~4.2--4.4.  Their terminal
quotient is identified with the dendritic quotient in Section~5, and
Theorem~5.23 makes it a dendrite.  Lemma~6.5 gives the induced convergence
action; minimality follows from equivariance and surjectivity.  Finally,
the fixed-point assertions are Lemmas~6.9 and~6.10.  All these constructions
apply to a continuum before any local connectivity is known.
\end{proof}

We seek a $\Gamma$-invariant full subset of $\Phi$ whose intervals are
finite. By part~\textup{(i)} of the proposition, it is already a
simplicial tree; a dense orbit of cut points will then show that it
contains all of $P$. To find it, we first use the splitting obstruction
to understand the final quotient $D(B)$ and the stabilizers of cut
points. We will return to the successive collapses after establishing
these dynamical facts.

\begin{lemma}\label{pret:parabolic-cuts}
Under the assumptions of this section, $D(B)$ is a point and every
$p\in P$ has infinite stabilizer containing no loxodromic element.
In particular, each such stabilizer contains a parabolic element of
infinite order fixing $p$.
\end{lemma}

\begin{proof}
Bowditch's dendrite obstruction
\cite[Lemma~1.4]{Bowditch1999Connected} applies to a finitely presented
one-ended group with no infinite torsion subgroup acting minimally by
convergence on a dendrite, provided every two-ended subgroup over which
the group splits is parabolic on that dendrite.  Here there are no such
splittings, so the last condition is vacuous; torsion-freeness implies
the torsion condition.  Proposition~\ref{pret:package} therefore gives
$D(B)$ a single point.

By \cite[Theorem~6.1]{Bowditch1999Treelike}, a minimal convergence action
of a finitely generated one-ended group with no infinite torsion subgroup
on a continuum has a nontrivial dendritic quotient whenever there is a
cut point which is not parabolic.
Its contrapositive shows that every point of $P$ is parabolic.  Here
``parabolic point'' means that its stabilizer is infinite and contains no
loxodromic element.  Any nonidentity element of this stabilizer has
infinite order because $\Gamma$ is torsion-free; the convergence-action
classification then makes it parabolic.
\end{proof}

Theorem~6.1 in the cited paper does not itself produce an invariant
finite-interval class when all cut points are parabolic.  We now prove
that assertion. The transfinite argument adapts the limit and
higher-successor steps of Bowditch's proof of Theorem~6.1, following
his Lemma~6.13. The parabolic fixed-point observation below permits
those steps in the present all-parabolic case.
The parabolic fixed-point description is recorded in
the remark after Lemma~6.10 of \cite{Bowditch1999Treelike}; we include
its short proof to make the extension explicit.

\begin{lemma}[Parabolic fixed points in the completion]\label{pret:parabolic-star}
Let $\gamma$ be a parabolic element with unique fixed point $a\in P$.
Every point of $\Phi$ fixed by $\gamma$ is either $a$ or adjacent to $a$.
Consequently, if an infinite-order element fixes two points of $\Phi$
whose interval is infinite, the element is loxodromic and both points
are terminal.
\end{lemma}

\begin{proof}
Every nonempty compact $\gamma$-invariant subset of $B$ contains $a$,
because the iterates of each point different from $a$ converge to $a$.
If $q\in\Phi$ is fixed, $R_B(q)$ is such a subset.  For $q\in P$,
uniqueness of the fixed point in $B$ gives $q=a$.  If $q\notin P$,
then $a\in R_B(q)$ and $\phi(a)=a\ne q$.  Formula~\eqref{pret:support}
therefore gives $a\in\Lambda(q)$, as required.

Any two points equal or adjacent to $a$ have a finite interval between
them, contained in the union of their intervals to $a$.  Thus a parabolic
element fixing a point of $P$ cannot have the two fixed points in the
last assertion.  A parabolic element whose fixed point is outside $P$
cannot have them either, by Proposition~\ref{pret:package}\textup{(iv)}.
The infinite-order classification and the loxodromic assertion in that
proposition complete the proof.
\end{proof}

\subsection{An invariant finite-interval class}

The final quotient $D(B)$ is a point. We now show that an invariant
class already appears at the first collapse, where its intervals are
finite in the original pretree. This is the stage that yields the
simplicial tree we need.

An equivalence relation on $\Phi$ is \emph{full} if all its classes are full.
For a full relation $\sim$, the quotient pretree has the induced
betweenness: a class $V$ lies between distinct classes $U,W$ when some
$v\in V$ lies between every $u\in U$ and $w\in W$.
Write $\sim_0$ for equality and define
\begin{equation}\label{pret:relations}
 \begin{aligned}
 x\sim_{\beta+1}y
 &\quad\Longleftrightarrow\quad
 [x,y]\text{ meets finitely many }\sim_\beta\text{-classes},\\
 \sim_\lambda&=\bigcup_{\beta<\lambda}\sim_\beta
 \qquad(\lambda\text{ a limit ordinal}).
 \end{aligned}
\end{equation}
These are the increasing invariant full relations in
Proposition~\ref{pret:package}\textup{(iii)}; their quotients remain
complete median pretrees.  In particular,
$x\sim_1y$ means precisely that $[x,y]$ is finite. Our goal is a fixed point of $\Phi/\sim_1$.

We record an elementary consequence of completeness that will identify
the direction of an edge between collapsed classes.  If disjoint
nonempty full subsets $U,V$ of a complete median pretree have full union,
a \emph{facing point of $V$ toward $U$} is a point $p\in V$ satisfying
\begin{equation}\label{pret:facing}
 [u,v]\cap V=[p,v]\qquad(u\in U, v\in V).
\end{equation}
There is at most one such point on each side, and there is one on at
least one side.  To see existence, choose $u\in U,v\in V$ and partition
$[u,v]$ by the two full sets.  Its initial part $A=[u,v]\cap U$ has
the form $[u,t]$ or $[u,t)$ by completeness.  In the closed case
$t\in U$.  In the half-open case put $m=\operatorname{med}(u,v,t)$.
If $m\ne t$, then $A=[u,m]$, reducing to the closed case.
Otherwise $t\in[u,v]$ and $t\notin A$, so $t\in V$.
Thus the initial part has an endpoint in $[u,v]$.  If the endpoint
belongs to $U$, it is the facing point there; otherwise it is the
facing point in $V$.  Interchanging $U,V$ if needed, denote this point
by $p\in V$.
The median of $u,u',p$ shows that the same endpoint works for every
$u'\in U$; the median of $u,v',p$ then gives
\eqref{pret:facing} for every $v'\in V$.  If $p,p'\in V$ both worked,
$\operatorname{med}(u,p,p')$ would force $p=p'$.  This is also
\cite[Lemma~6.13]{Bowditch1999Treelike}.

\begin{lemma}\label{pret:invariant-class}
Under the assumptions of this section, $\Gamma$ preserves a
finite-interval class of $\Phi$.
\end{lemma}

\begin{proof}
Lemma~\ref{pret:parabolic-cuts} says that the terminal quotient in
\eqref{pret:relations} is a point.  Let $\alpha$ be the least ordinal
for which $\Phi/\sim_\alpha$ has a point fixed by $\Gamma$.
Proposition~\ref{pret:package}\textup{(ii)} gives $\alpha\ne0$.
We shall rule out limits and successors greater than one.

\emph{A limit cannot be first.}
Suppose $\alpha$ is a limit ordinal, and choose $x$ in an invariant
$\sim_\alpha$-class.  For each member $s$ of a finite generating set,
$x\sim_\alpha sx$ holds at some stage $\beta_s<\alpha$.
Let $\beta$ be the maximum of these finitely many stages.  Invariance
and transitivity of $\sim_\beta$ give $x\sim_\beta gx$ for every word
$g$ in the generators and their inverses.  Thus the $\sim_\beta$-class
of $x$ is invariant, contrary to the choice of $\alpha$.

\emph{A higher successor creates a simplicial tree.}
Suppose $\alpha=\beta+1$ with $\beta\ge1$, and let $\Xi\subset\Phi$
be an invariant $\sim_\alpha$-class.  Its $\sim_\beta$-classes form a
simplicial tree $\Sigma=\Xi/\sim_\beta$.  The action on this tree has
no fixed vertex, by the minimality of $\alpha$.

For adjacent classes $U,V\in\Sigma$, their union is full in $\Phi$.
The facing-point observation supplies a point on at least one side.
There cannot be facing points on both sides: they would be adjacent in
$\Phi$, hence $\sim_1$-equivalent and therefore $\sim_\beta$-equivalent.
This contradicts $U\ne V$.  Exactly one side has a facing point, so an
automorphism cannot interchange $U$ and $V$.  The action on $\Sigma$
has no edge inversions.

We next show that every edge of $\Sigma$ with infinite stabilizer is a
leaf edge.  This is the only step in which the dynamics on $B$ enters
the ordinal argument.

Let $UV$ have infinite stabilizer, and write $p\in V$ for its facing
point.  Choose a nonidentity stabilizing element $\gamma$; it has infinite
order.  Since the edge is not inverted, $\gamma$ preserves $U,V$ and
fixes their uniquely specified point $p$.
We claim that $\gamma$ fixes another point outside $V$.
For $y\in U$ put
\[
 x=\operatorname{med}(p,y,\gamma y).
\]
Fullness gives $x\in[y,\gamma y]\subset U$, and both $x,\gamma x$
lie in the linearly ordered interval $[p,\gamma y]$.
If $x=\gamma x$, the claim is proved.  Otherwise, after replacing
$(\gamma,x)$ by $(\gamma^{-1},\gamma x)$ if needed, we have
$x\in(p,\gamma x)$.
The intervals $[p,\gamma^n x]$, $n\ge0$, are then strictly increasing.
Their union is a full directed arc with first point $p$ and no last point.
Completeness supplies an endpoint $z$ with
\begin{equation}\label{pret:endpoint}
 \bigcup_{n\ge0}[p,\gamma^n x]=[p,z).
\end{equation}
The endpoint is unique in a median pretree
\cite[Lemma~3.13 and the following observation]{Bowditch1999Treelike}.
The union in \eqref{pret:endpoint} is $\gamma$-invariant, so $z$ is fixed.
Since $x\in(p,z)$ but $x\notin V$, fullness of $V$ implies $z\notin V$.
This proves the claim.

\begin{figure}[!htbp]
\centering
\begin{tikzpicture}[x=0.95cm,y=0.8cm,>=Stealth]
 \fill[gray!10] (-0.9,-0.6) rectangle (0,0.7);
 \node at (-0.55,0.35) {$V$};
 \draw[thick] (0,0)--(5.9,0);
 \fill (0,0) circle (2pt) node[below=5pt] {$p$};
 \fill (1.4,0) circle (2pt) node[below=5pt] {$x$};
 \fill (2.65,0) circle (2pt) node[below=5pt] {$\gamma x$};
 \fill (3.65,0) circle (2pt) node[below=5pt] {$\gamma^2x$};
 \node at (4.6,0.2) {$\cdots$};
 \fill (5.9,0) circle (2pt) node[below=5pt] {$z$};
 \draw[->] (1.4,0.65)--(3.9,0.65);
 \node[above] at (2.7,0.65) {$U$};
 \node[right] at (6.2,0) {$z\notin V$};
\end{tikzpicture}
\caption{The interval argument for an edge between higher-stage classes.
The element fixes the facing point $p$.  If it moves points of $U$
away from $p$, completeness supplies a second fixed endpoint $z$ outside
$V$.  The picture depicts betweenness in $\Phi$, not an arc in $B$.}
\label{fig:pretree-endpoint}
\end{figure}

Let $q\notin V$ be the second fixed point.  The points $p,q$ belong
to different $\sim_\beta$-classes; as $\beta\ge1$, their original
interval is infinite.  Lemma~\ref{pret:parabolic-star} makes $p$ terminal
in $\Phi$.  For every $v\in V$ and $u\in U$, formula~\eqref{pret:facing}
puts $p\in[u,v]$.  Since $p\ne u$, terminality forces $v=p$.
Thus $V=\{p\}$.  This vertex is a leaf of $\Sigma$: if it had two
neighbors, $p$ would lie strictly between representatives of them in
$\Phi$, again contradicting terminality.

Delete every edge with infinite stabilizer, together with its leaf
endpoint.  The remaining vertices form a nonempty invariant subtree.
Indeed, removing selected leaves preserves all paths between remaining
vertices.  If nothing remained, the original tree would have at most
two vertices.  Such an action without inversions has a fixed vertex,
which has already been excluded.  All remaining edges have finite
stabilizer, and a fixed vertex of the remaining tree would also be fixed
in $\Sigma$.

A finitely generated one-ended group acting without inversions on a
simplicial tree with finite edge stabilizers fixes a vertex
\cite[Lemma~6.6]{Bowditch1999Treelike}.  Applying this fact gives the
desired contradiction.

The least ordinal is therefore $\alpha=1$, which is exactly the
assertion of the lemma.
\end{proof}

\Needspace{14\baselineskip}
\subsection{The simplicial tree spanning the cut set}

The ordinal argument is complete.  We now work only with an ordinary
simplicial tree, whose vertices retain the separation information in $B$.

\begin{proposition}\label{pret:tree}
Let $\Gamma$ be a finitely presented torsion-free one-ended group with no
nontrivial splitting over a finite or two-ended subgroup.  Suppose it acts
minimally by convergence on a nondegenerate compact metric continuum $B$
whose set $P$ of global cut points is nonempty.  Then there is a
nontrivial minimal simplicial $\Gamma$-tree $\mathcal T$, with no edge
inversions and finite quotient, such that:
\begin{enumerate}[label=\textup{(\roman*)}]
\item $P$ is an invariant subset of its vertices, and every edge has one
endpoint in $P$ and one outside $P$;
\item $\mathcal T$ is spanned by $P$; its other vertices have degree at
least two;
\item every edge stabilizer is infinite;
\item if $a,b,q\in P$ and $q$ separates $a,b$ in $B$, then
$q\in[a,b]_{\mathcal T}$.
\end{enumerate}
Every point of $P$ is fixed by an infinite-order parabolic element.
\end{proposition}

\begin{proof}
Let $C\subset\Phi$ be the invariant finite-interval class supplied by
Lemma~\ref{pret:invariant-class}.  It contains more than one point,
because $\Phi$ has no fixed point.  Therefore its simplicial tree has
an edge, and Proposition~\ref{pret:package}\textup{(i)} places one
endpoint of that edge in $P$.

Choose $p\in C\cap P$.  The orbit $\Gamma p$ is dense in $B$ and lies
in $C$.  If $q\in P$, take an open separation of $B\setminus\{q\}$.
Both sides are nonempty open subsets of $B$, so each meets $\Gamma p$.
For orbit points $a,b$ on opposite sides, separation gives
$q\in[a,b]_{\Phi}$.  Since $C$ is full, $q\in C$.  We have proved
$P\subset C$.

Let $\mathcal T$ be the subtree of $C$ spanned by $P$, that is, the union
of the intervals between its points.  Its intervals agree with those of
$\Phi$, proving the asserted separation compatibility.  The adjacency
property gives the indicated bipartition and rules out edge inversions.
A vertex outside $P$ cannot be a leaf of a tree spanned by $P$.

To prove minimality, let $\mathcal S\subset\mathcal T$ be a nonempty
invariant subtree.  A singleton would be a fixed point of $\Phi$.
Otherwise $\mathcal S$ contains an edge, hence a point of $P$.
Its orbit is again dense in $B$; the separation argument of the preceding
paragraph gives $P\subset\mathcal S$.  Thus
$\mathcal S=\mathcal T$.

Choose a vertex $o$ and a finite generating set $S$ for $\Gamma$.
The union of the finitely many finite paths $[o,so]$, $s\in S$, is a
finite subtree after adjoining $o$.  Its translates form a connected
invariant subtree and therefore cover $\mathcal T$ by minimality.
Consequently the quotient has finitely many vertex and edge orbits.

Finally suppose an edge had finite stabilizer.  Collapsing the edges
outside its orbit gives a nontrivial tree with one edge orbit and a
minimal action without inversions.  Bass--Serre theory then yields a
nontrivial splitting over that edge stabilizer, contrary to the
hypotheses.  Thus every edge stabilizer is infinite.  The parabolic
assertion is Lemma~\ref{pret:parabolic-cuts}.
\end{proof}
